\documentclass[10pt,a4paper]{amsart}
\usepackage[margin=19mm]{geometry}
\usepackage{amsmath,amssymb,mathtools}
\usepackage[scr=rsfs,cal=euler]{mathalfa}
\usepackage{xfrac}
\usepackage[unicode,hidelinks,bookmarks=false]{hyperref}
\newcommand{\ie}{i.e.\ }
\newcommand{\cf}{cf.\ }
\newcommand{\resp}{respectively\ }
\newcommand{\Subsection}{\subsection}
\providecommand{\nobreakdash}{\nobreak\hbox{-}}
\setlength{\emergencystretch}{2em}
\multlinegap0pt
\usepackage{mathtools}
\usepackage{amssymb}
\usepackage{bm}
\usepackage{algorithm}\usepackage{algpseudocode}
\usepackage[shortlabels]{enumitem}
\usepackage{xcolor}
\newtheorem{lemma}{Lemma}
\usepackage{cleveref}
\crefname{lemma}{Lemma}{Lemmas}
\newcommand{\R}{\mathbb{R}}
\DeclareMathOperator{\St}{St}
\newcommand{\bze}{\bm{0}}
\DeclareMathOperator{\col}{col}
\newcommand{\opn}[1]{\left\lVert #1\right\rVert_{\mathrm{op}}}
\newcommand{\trans}{\top}
\title{Bias-Corrected Subspace Intersection: Minimax-Optimal Shared Subspace Estimation in Multi-View Data}
\author{Xianwen Song, Yuepeng Yang, Cong Ma}
\date{Source: arXiv:2609.05617v1 (CC BY 4.0)}
\begin{document}
\maketitle

\noindent\emph{Source and licence.} Bias-Corrected Subspace Intersection: Minimax-Optimal Shared Subspace Estimation in Multi-View Data. Version arXiv:2609.05617v1 (CC BY 4.0), distributed by arXiv under the Creative Commons Attribution 4.0 International licence, http://creativecommons.org/licenses/by/4.0/, as recorded in the arXiv licence metadata for this exact version and shown on the abstract page https://arxiv.org/abs/2609.05617v1. Full licence notice: \texttt{licenses/arxiv-2609.05617-CC-BY-4.0.txt}.\par\smallskip

\noindent\textit{Editorial scope.} Counted unit: the article's lemma lem:nullspace (source0.tex lines 1521--1541). The article's complete original proof of it is delivered with it (source0.tex lines 2364--2530, one \texttt{proof} environment of 167 lines, which the article places away from the statement). No mathematics is written, repaired or re-derived: the statement and the proof are the article's own lines, in the article's own notation, and no retained line is substituted or re-flowed.  No further source-local context is needed: every cross-reference of every retained line resolves inside the two delivered spans.  Nothing was omitted from any retained span, so the package carries no omission record.  No retained line contains a citation, so the wrapper carries no bibliography and no bibliography entry.  No retained line contains a figure environment, an image-inclusion command or drawing code, so no image is delivered and the package declares no asset dependency.  Editorial formatting only, in addition: an AMS \texttt{amsart} preamble that restates the article's own declarations and macros used by the retained lines, namely the environments \texttt{lemma} on separate counters, together with the standard packages those lines need.  The retained environments print in this document as Lemma 1 in order of appearance, whereas the article prints the same items with the numbering of its own full document.  The complete native source of the article as distributed in the arXiv e-print for this exact version is bundled unchanged as \texttt{sources/209451/source0.tex}.
\begin{lemma}
\label{lem:nullspace}
Assume that $\dim\ker(\bm G_0)=r$ for some $1\le r< d$, and $\bm G = \bm G_0 + \bm \Delta$ with
\[
  \opn{\bm\Delta_{++}}\le\frac14,
  \qquad
  \varepsilon_\Delta
  \le\frac{1}{16}\lambda_{\min}^{+}(\bm G_0).
\]
Then $\bm G$ has exactly $r$ eigenvalues in
$[-\varepsilon_\Delta,\varepsilon_\Delta]$,
while all remaining eigenvalues exceed $4\varepsilon_\Delta$.
If $\bm Q\in\operatorname{St}(d,r)$ spans the invariant subspace
associated with these $r$ eigenvalues, then
\begin{equation}
\label{eq:abstract-relative-nullspace-conclusion}
  \opn{\bm G_0^{1/2}\bm Q}
  \le
  2\opn{\bm\Delta_{+0}}.
\end{equation}
\end{lemma}

\begin{proof}[Proof of \cref{lem:nullspace}]
The proof has three parts.  We first obtain a lower bound for the entire
spectrum of $\bm G$.  We then use the min--max principle to separate the
$r$ eigenvalues near zero from the remaining spectrum.  Finally, we project
the invariant-subspace equation onto $\ker(\bm G_0)^\perp$ to control the
relative residual.


Choose orthonormal bases
$\bm Q_0\in\St(d,r)$ and
$\bm Q_+\in\St(d,d-r)$
for $\ker(\bm G_0)$ and $\ker(\bm G_0)^\perp$, respectively, and define the
positive spectral block
$\bm K\coloneqq\bm Q_+^\trans\bm G_0\bm Q_+$.
Then $\bm K\succ\bze$ and
$\lambda_{\min}(\bm K)=\lambda_{\min}^{+}(\bm G_0)$.
We identify the relative perturbation operators with their nonzero coordinate
blocks,
\[
  \bm\Delta_{00}
  =
  \bm Q_0^\trans\bm\Delta\bm Q_0,
  \qquad
  \bm\Delta_{+0}
  =
  \bm K^{-1/2}\bm Q_+^\trans\bm\Delta\bm Q_0,
  \qquad
  \bm\Delta_{++}
  =
  \bm K^{-1/2}
  \bm Q_+^\trans\bm\Delta\bm Q_+
  \bm K^{-1/2}.
\]
This identification preserves their operator norms.  The assumptions of
\cref{lem:nullspace} imply
\begin{equation}
\label{eq:null-smallness}
  \opn{\bm\Delta_{++}}\le\frac14,
  \qquad
  \varepsilon_\Delta
  \le\frac{1}{16}\lambda_{\min}^{+}(\bm G_0).
\end{equation}

\paragraph{Lower bound for the spectrum.}
Write any $\bm w\in\R^d$ uniquely as
$\bm w=\bm Q_0\bm u+\bm Q_+\bm v$, and set
$\bm t\coloneqq\bm K^{1/2}\bm v$.  In these coordinates,
\begin{equation}
\label{eq:null-quadratic-form}
  \bm w^\trans\bm G\bm w
  =
  \|\bm t\|_2^2
  +
  \bm t^\trans\bm\Delta_{++}\bm t
  +
  2\bm t^\trans\bm\Delta_{+0}\bm u
  +
  \bm u^\trans\bm\Delta_{00}\bm u.
\end{equation}
By \cref{eq:null-smallness}, the first two terms are at least
$\frac12\|\bm t\|_2^2$.  Young's inequality gives
$2|\bm t^\trans\bm\Delta_{+0}\bm u|
\le\frac12\|\bm t\|_2^2+
2\|\bm\Delta_{+0}\|_{\mathrm{op}}^2\|\bm u\|_2^2$.
Consequently,
\[
  \bm w^\trans\bm G\bm w
  \ge
  -\left(
    \opn{\bm\Delta_{00}}
    +
    2\opn{\bm\Delta_{+0}}^2
  \right)
  \|\bm u\|_2^2
  \ge
  -\varepsilon_\Delta\|\bm w\|_2^2.
\]
It follows that $\lambda_{\min}(\bm G)\ge-\varepsilon_\Delta$.

\paragraph{Separation of the near-zero spectrum.}
Order the eigenvalues as
$\lambda_1(\bm G)\ge\cdots\ge\lambda_d(\bm G)$.
On $\ker(\bm G_0)=\col(\bm Q_0)$, the quadratic form of $\bm G$ is
$\bm u^\trans\bm\Delta_{00}\bm u$.  Courant--Fischer therefore gives
\[
  \lambda_{d-r+1}(\bm G)
  \le
  \opn{\bm\Delta_{00}}
  \le
  \varepsilon_\Delta.
\]
Thus at least $r$ eigenvalues are at most $\varepsilon_\Delta$.

To control the remaining spectrum, apply Courant--Fischer on
$\ker(\bm G_0)^\perp=\col(\bm Q_+)$.  In this case $\bm u=\bze$, and
\cref{eq:null-quadratic-form} gives
\[
  \bm w^\trans\bm G\bm w
  \ge
  \lambda_{\min}^{+}(\bm G_0)
  \bigl(
    1-\opn{\bm\Delta_{++}}
  \bigr)
  \|\bm w\|_2^2.
\]
By \cref{eq:null-smallness}, the coefficient on the right is at least
$12\varepsilon_\Delta$, which is strictly larger than
$4\varepsilon_\Delta$.  Hence
$\lambda_{d-r}(\bm G)>4\varepsilon_\Delta$.
Combining this separation with the preceding lower bound shows that exactly
$r$ eigenvalues lie in $[-\varepsilon_\Delta,\varepsilon_\Delta]$, while all remaining eigenvalues
exceed $4\varepsilon_\Delta$.

\paragraph{Bound for $\bm G_0^{1/2}\bm Q$.}
Let $\bm Q\in\St(d,r)$ span the invariant subspace associated with
the $r$ eigenvalues in $[-\varepsilon_\Delta,\varepsilon_\Delta]$, and set
$\bm\Lambda\coloneqq\bm Q^\trans\bm G\bm Q$.
Then $\bm G\bm Q=\bm Q\bm\Lambda$ and
$\|\bm\Lambda\|_{\mathrm{op}}\le \varepsilon_\Delta$.

To express the desired residual in the positive spectral coordinates, define
$\bm H\coloneqq\bm K^{1/2}\bm Q_+^\trans\bm Q$.
Projecting the invariant-subspace equation onto $\col(\bm Q_+)$ and
multiplying by $\bm K^{-1/2}$ gives
\begin{equation}
\label{eq:nullspace-H-equation}
  \bm H
  =
  -\bm\Delta_{+0}\bm Q_0^\trans\bm Q
  -\bm\Delta_{++}\bm H
  +\bm K^{-1}\bm H\bm\Lambda.
\end{equation}
Taking operator norms, using
$\|\bm Q_0^\trans\bm Q\|_{\mathrm{op}}\le1$, and recalling that
$\|\bm K^{-1}\|_{\mathrm{op}}
=1/\lambda_{\min}^{+}(\bm G_0)$, we obtain
\[
  \opn{\bm H}
  \le
  \opn{\bm\Delta_{+0}}
  +
  \left(
    \opn{\bm\Delta_{++}}
    +
    \frac{\opn{\bm\Lambda}}
         {\lambda_{\min}^{+}(\bm G_0)}
  \right)
  \opn{\bm H}.
\]
Since $\|\bm\Lambda\|_{\mathrm{op}}\le \varepsilon_\Delta$, the factor multiplying
$\|\bm H\|_{\mathrm{op}}$ is at most $1/2$ by
\cref{eq:null-smallness}.  Therefore
$\|\bm H\|_{\mathrm{op}}
\le2\|\bm\Delta_{+0}\|_{\mathrm{op}}$.
Finally,
\[
  \opn{\bm G_0^{1/2}\bm Q}
  =
  \opn{\bm K^{1/2}\bm Q_+^\trans\bm Q}
  =
  \opn{\bm H}
  \le
  2\opn{\bm\Delta_{+0}}.
\]
This is exactly the bound in
\cref{eq:abstract-relative-nullspace-conclusion}, so the result follows.
\end{proof}

\end{document}
